\section{Conclusion: What LLM Advice Flattens}

Instead of treating alignment as an intrinsic property of LLMs---whether models are helpful, harmless, or honest---it is better understood as a situated relationship between model behavior and the social contexts in which it is interpreted and evaluated. A model aligned with the averaged norms of diffuse internet training is misaligned with any concentrated community's specific wisdom. A model trained to avoid prescriptive language is misaligned with contexts where adequate advice includes authorizing action. 

The question for researchers and designers is not whether LLMs are aligned, but who might benefit from that alignment, and who might suffer. For advice-seekers facing clear-cut situations of harm, the consequences of the current study are practical: the permission to act that communities grant and LLMs do not is not a stylistic difference but a functional one. This is consistent with understanding safety alignment as a specific form of algorithmic moral governance, whose practical effect is to prioritize corporate risk-aversion and therapeutic self-management over the situated, directive friction often necessary for escaping harm.

The concept of vernacular competence, and the method of reading AI outputs against concentrated community norms, offer tools for investigating these dynamics beyond the case examined in this article. Any domain where communities have developed situated practices of authorization possesses a vernacular competence that LLM contrast can render visible. What the analysis shows is that assistant-style AI carries a portable normative style that flattens situated modes of judgment, and that this matters most precisely where communities treat directive action as necessary rather than excessive.